% =============================================================================
% INTRODUCTION - CITA Paper in ALKALI Format
% =============================================================================

\vspace{-2.5em}
\section{Why Instruction-Based Alignment?}
\label{sec:introduction}
\vspace{-0.5em}

LLMs are increasingly deployed not as single-purpose chatbots, but as \textbf{agentic backbones} that sit behind many products and workflows---customer support, internal research assistants, compliance copilots, tutoring, code review, and creative ideation~\cite{brown2020language,openai2023gpt4,touvron2023llama}. In practice, these workflows differ less in \emph{capability} than in \emph{alignment posture}: the same underlying competence must be expressed with different safety thresholds, tone, risk tolerance, and disclosure norms depending on \textbf{role}, \textbf{context}, and \textbf{policy}. Put plainly, modern deployments require \textbf{one model, many behavioral contracts}.







% Requires: \usepackage{tabularx,booktabs,multirow,array,pifont}
% Icons (FontAwesome): \usepackage{fontawesome5}
% For compact cells: \newcolumntype{L}[1]{>{\raggedright\arraybackslash}p{#1}}
% Optional: \renewcommand{\arraystretch}{1.12}

% =======================
% Preamble (make sure)
% =======================
% \usepackage{tabularx,booktabs,array}
% \usepackage{xurl}
% \usepackage[hidelinks]{hyperref}
% \usepackage{fontawesome5} % if allowed; otherwise keep the SafeFA fallbacks below
% \usepackage{pifont}       % fallback icons
%
% % Robust icon fallback (prevents "Undefined control sequence")
% \makeatletter
% \newcommand{\SafeFA}[2]{\ifcsname #1\endcsname \csname #1\endcsname \else #2\fi}
% \makeatother
% \newcommand{\IcSupport}{\SafeFA{faHeadphones}{\ding{72}}}
% \newcommand{\IcResearch}{\SafeFA{faFlask}{\ding{46}}}
% \newcommand{\IcCompliance}{\SafeFA{faGavel}{\ding{115}}}
% \newcommand{\IcCreative}{\SafeFA{faLightbulb}{\ding{72}}}
% \newcommand{\IcChild}{\SafeFA{faChild}{\ding{43}}}
% \newcommand{\IcSecurity}{\SafeFA{faShieldAlt}{\ding{110}}}
%
% % Helps reduce overfull boxes in tight columns
% \setlength{\emergencystretch}{2em}

% Preamble helpful (optional but recommended)
% \usepackage{tabularx,booktabs,array}
% \usepackage{caption}
% \captionsetup[table]{justification=raggedright,singlelinecheck=false}

% Preamble (ensure):
% \usepackage{tabularx,booktabs,array}
% \usepackage{pifont}
% \usepackage{fontawesome5} % if allowed
% (and your \IcSupport ... icon macros)

% ============================
% One-column version (no table)
% NeurIPS-clean, compact, LaTeX-safe
% ============================
% Preamble requirements:
% \usepackage{pifont}   % for \ding
% \usepackage{xurl}
% \usepackage[hidelinks]{hyperref}

% ============================
% One-column rows with pastel BACKGROUNDS (safe, NeurIPS-clean)
% ============================
% Preamble requirements:
% \usepackage{xcolor}
% \usepackage{pifont}
% \usepackage{xurl}
% \usepackage[hidelinks]{hyperref}
%
% Optional but recommended (better spacing/linebreaks):
% \usepackage{microtype}

% ---- Palette (pastel, print-friendly) ----
\definecolor{Ink}{HTML}{111111}
\definecolor{Line}{HTML}{D7DEE7}
\definecolor{Teal}{HTML}{2F7F73}
\definecolor{Blue}{HTML}{2B6CB0}
\definecolor{Violet}{HTML}{6B46C1}
\definecolor{Amber}{HTML}{B7791F}
\definecolor{Rose}{HTML}{B83280}
\definecolor{Slate}{HTML}{374151}

\definecolor{TealBg}{HTML}{EAF6F3}
\definecolor{BlueBg}{HTML}{EEF5FF}
\definecolor{VioletBg}{HTML}{F3EEFF}
\definecolor{AmberBg}{HTML}{FFF6E5}
\definecolor{RoseBg}{HTML}{FCEFF7}
\definecolor{SlateBg}{HTML}{F2F4F7}

% ---- Helpers ----
\newcommand{\K}[1]{\textbf{\textcolor{Ink}{#1}}}
\newcommand{\EM}[1]{\emph{#1}}
\newcommand{\SEP}{\;\textcolor{Line}{\textbar}\;}

% Icons (pifont-safe)
\newcommand{\Icsup}{\ding{72}}
\newcommand{\Icre}{\ding{46}}
\newcommand{\Iccm}{\ding{115}}
\newcommand{\Icbulb}{\ding{72}}
\newcommand{\Icchild}{\ding{43}}
\newcommand{\Icsec}{\ding{110}}

% Colored labels
\newcommand{\Lrisk}{\textbf{\textcolor{Rose}{Risk:}}}
\newcommand{\Lout}{\textbf{\textcolor{Blue}{Out:}}}
\newcommand{\Ltrace}{\textbf{\textcolor{Violet}{Trace:}}}
\newcommand{\Ltone}{\textbf{\textcolor{Teal}{Tone:}}}
\newcommand{\Lstyle}{\textbf{\textcolor{Teal}{Style:}}}
\newcommand{\Lcontent}{\textbf{\textcolor{Teal}{Content:}}}
\newcommand{\Lref}{\textbf{\textcolor{Amber}{Refusal:}}}
\newcommand{\Lcav}{\textbf{\textcolor{Violet}{Caveat:}}}
\newcommand{\Lex}{\textbf{\textcolor{Slate}{Ex:}}}
\newcommand{\Lpr}{\EM{Prompt:}}
\newcommand{\Lexp}{\EM{Expected:}}
\newcommand{\Lnot}{\EM{Not:}}

% Thin separators
\newcommand{\ThinRule}{\noindent\textcolor{Line}{\rule{\linewidth}{1pt}}}

% ---- Row background wrapper (NO extra packages; uses \colorbox + \parbox) ----
% Usage: \RowBox{<bgcolor>}{<content>}
\newcommand{\RowBox}[2]{%
  \noindent\colorbox{#1}{%
    \parbox{\dimexpr\linewidth-2\fboxsep\relax}{#2}%
  }%
}

\vspace{-1em}


%\vspace{0.6mm}
\ThinRule
%\vspace{0.7mm}

% ============================
% Row 1
% ============================
\RowBox{TealBg}{%
  %\vspace{0.7mm}
  \noindent\K{\textcolor{Teal}{\Icsup}\ \ Customer Support}
  \hfill {\footnotesize\textsc{SAFE--EMPATHETIC}}\\[-0.2mm]
  {\scriptsize \Lrisk\ \textsc{CONSERVATIVE} \SEP \Lout\ checklist + escalate \SEP \Lref\ help}\\[-0.2mm]
  {\scriptsize \Ltone\ warm, de-escalating \SEP \Lstyle\ minimal jargon \SEP \K{Disclosure:} avoid overconfidence}\\[-0.2mm]
  {\scriptsize \Lex\ \Lpr\ ``My account was hacked---what do I do?'' \SEP \Lexp\ safe checklist + escalation \SEP \Lnot\ risky security steps}
  %\vspace{0.7mm}
}

%\vspace{0.6mm}
\ThinRule
%\vspace{0.7mm}

% ============================
% Row 2
% ============================
\RowBox{BlueBg}{%
  %\vspace{0.7mm}
  \noindent\K{\textcolor{Blue}{\Icre}\ \ Internal Research}
  \hfill {\footnotesize\textsc{ACCURATE--NUANCED}}\\[-0.2mm]
  {\scriptsize \Lrisk\ \textsc{BALANCED} \SEP \Lout\ cite + caveat \SEP \Lref\ only if disallowed}\\[-0.2mm]
  {\scriptsize \Ltone\ professional \SEP \K{Epistemics:} cite sources / uncertainty \SEP \Lcontent\ allow nuance}\\[-0.2mm]
  {\scriptsize \Lex\ \Lpr\ ``Summarize failure modes of DPO and RLHF.'' \SEP \Lexp\ structured comparison + citations \SEP \Lnot\ blanket refusal boilerplate}
  %\vspace{0.7mm}
}

%\vspace{0.6mm}
\ThinRule
%\vspace{0.7mm}

% ============================
% Row 3
% ============================
\RowBox{VioletBg}{%
  %\vspace{0.7mm}
  \noindent\K{\textcolor{Violet}{\Iccm}\ \ Compliance Review}
  \hfill {\footnotesize\textsc{POLICY--STRICT}}\\[-0.2mm]
  {\scriptsize \Lrisk\ \textsc{HIGH} \SEP \Ltrace\ log rationale \SEP \Lout\ conservative \SEP \Lref\ firm}\\[-0.2mm]
  {\scriptsize \Ltone\ formal \SEP \K{Action:} flag violations \SEP \K{Outputs:} conservative summaries}\\[-0.2mm]
  {\scriptsize \Lex\ \Lpr\ ``Draft an email to request medical records.'' \SEP \Lexp\ privacy-aware template + consent \SEP \Lnot\ unnecessary sensitive details}
  %\vspace{0.7mm}
}

%\vspace{0.6mm}
\ThinRule
%\vspace{0.7mm}

% ============================
% Row 4
% ============================
\RowBox{AmberBg}{%
  %\vspace{0.7mm}
  \noindent\K{\textcolor{Amber}{\Icbulb}\ \ Creative}
  \hfill {\footnotesize\textsc{EXPLORATORY--BOUNDED}}\\[-0.2mm]
  {\scriptsize \Lrisk\ \textsc{PERMISSIVE} \SEP \Lout\ high-diversity \SEP \Lcav\ label speculation}\\[-0.2mm]
  {\scriptsize \Ltone\ playful \SEP \Lstyle\ divergent ideas \SEP \K{Safety:} avoid disallowed content}\\[-0.2mm]
  {\scriptsize \Lex\ \Lpr\ ``Give 20 wild taglines for a new AI tutor.'' \SEP \Lexp\ high-diversity options \SEP \Lnot\ unnecessary refusal}
  %\vspace{0.7mm}
}

%\vspace{0.6mm}
\ThinRule
%\vspace{0.7mm}

% ============================
% Row 5
% ============================
\RowBox{RoseBg}{%
  %\vspace{0.7mm}
  \noindent\K{\textcolor{Rose}{\Icchild}\ \ Child Education}
  \hfill {\footnotesize\textsc{AGE--APPROPRIATE}}\\[-0.2mm]
  {\scriptsize \Lrisk\ \textsc{STRICT} \SEP \Lout\ scaffold + examples \SEP \Lref\ safe redirect}\\[-0.2mm]
  {\scriptsize \Ltone\ encouraging \SEP \Lcontent\ no mature topics \SEP \K{Pedagogy:} scaffolding + examples}\\[-0.2mm]
  {\scriptsize \Lex\ \Lpr\ ``Explain how the Internet works.'' \SEP \Lexp\ simple analogies \SEP \Lnot\ overly technical/unsafe detail}
  %\vspace{0.7mm}
}

%\vspace{0.6mm}
\ThinRule
%\vspace{0.7mm}

% ============================
% Row 6
% ============================
\RowBox{SlateBg}{%
  %\vspace{0.7mm}
  \noindent\K{\textcolor{Slate}{\Icsec}\ \ Security}
  \hfill {\footnotesize\textsc{DEFENSIVE--RESPONSIBLE}}\\[-0.2mm]
  {\scriptsize \Lrisk\ \textsc{STRICT} \SEP \Lout\ defensive only \SEP \Lref\ block misuse}\\[-0.2mm]
  {\scriptsize \Ltone\ calm \SEP \Lcontent\ defensive guidance only \SEP \K{Action:} detection/mitigation}\\[-0.2mm]
  {\scriptsize \Lex\ \Lpr\ ``How do I test if our API is vulnerable to injection?'' \SEP \Lexp\ authorized testing checklist \SEP \Lnot\ exploit/intrusion steps}
  %\vspace{0.7mm}
}

\vspace{0.4mm}
\ThinRule
\\
\noindent\K{One backbone, many alignment contracts.}
Instruction-driven alignment enables explicit, auditable \K{policy switching} across agent roles without proliferating checkpoints.
%\vspace{-1em}






\textbf{Static alignment is mismatched to runtime reality.}
Today’s alignment pipelines---RLHF~\cite{ouyang2022training,stiennon2020learning}, DPO~\cite{rafailov2023direct}, and variants~\cite{azar2023general,ethayarajh2024kto}---largely produce a \textbf{single, fixed policy per checkpoint}.
Once trained, the model’s ``alignment stance'' becomes a property of the weights: it tends to respond \emph{the same way} whether it is serving a child, a clinician, a security researcher, or a creative writer.
This rigidity forces a deployment trade-off that is now commonplace:
either (i) maintain \textbf{multiple separately-aligned checkpoints} (expensive to train, validate, version, and govern), or (ii) accept \textbf{one-size-fits-all behavior} (often suboptimal, sometimes unsafe).
As organizations adopt agent platforms, this trade-off becomes a bottleneck---not only in cost, but in \textbf{governance velocity}: policies evolve faster than \textbf{retraining cycles}.

 

\textbf{Alignment should be a controllable interface, not a frozen artifact.}
Agent orchestrators already rely on system prompts to route tasks, enforce formats, and set tool-use constraints.
We argue that alignment should be treated similarly: a model should accept \textbf{natural-language alignment instructions} that specify the desired behavioral contract \emph{at inference time}.
This turns alignment into a \textbf{runtime control channel}---auditable, updateable, and role-aware---rather than a training-only property.
The goal is not prompt ``hacks'' or brittle jailbreak-like steering, but \textbf{learned, stable switching} where the same prompt elicits different aligned behaviors \emph{because the policy instruction changes}, not because the user prompt is.
rewritten. 


{\scriptsize
%\vspace{-0.5em}
\begin{tcolorbox}[
    colback=Teal!5!White,
    colframe=Teal!55!black,
    title=\textbf{Research Questions},
    fonttitle=\bfseries,
    boxrule=0.5pt,
    arc=2mm,
    left=2mm, right=2mm, top=1mm, bottom=1mm
]
\begin{itemize}[leftmargin=*, itemsep=2pt, topsep=2pt]

    \item[\ding{227}] \textbf{RQ1: Can alignment policies be switched at inference time?}\\
    \textit{Can a single model reliably exhibit different alignment behaviors conditioned on natural-language alignment instructions, \emph{without} retraining or maintaining multiple checkpoints?}

    \item[\ding{227}] \textbf{RQ2: What learning objective enables \emph{stable} instruction-conditioned alignment?}\\
    \textit{How must preference optimization be modified to couple alignment instructions with behavioral change, and what role does an explicit KL anchor play in preventing collapse while preserving switchability?}

    \item[\ding{227}] \textbf{RQ3: How do we measure \emph{instruction-alignment} vs.\ instruction-following?}\\
    \textit{What controlled evaluations can separate deep, robust policy embodiment from superficial instruction compliance (e.g., format/tone adherence)?}

    \item[\ding{227}] \textbf{RQ4: Does switchability generalize beyond curated prompts and survive safety--utility trade-offs?}\\
    \textit{Across heterogeneous benchmarks (truthfulness, conditional safety, controllable length, and alignment quality), does instruction-driven switching remain consistent, predictable, and non-degrading?}

\end{itemize}
\end{tcolorbox}
}
\vspace{-1em}


%\subsection{Motivating Scenario: One Backbone, Many Agents}
%\label{sec:motivation}



% =============================================================================
% TEASER FIGURE: Real responses from all 5 trainers (Google Sheets screenshot)
% =============================================================================
\begin{figure*}[ht!]
\centering
\includegraphics[width=0.95\textwidth, trim=0 225 0 0, clip]{figures/evaluation/combined_plots/teaser_single_prompt_idx287.pdf}
\vspace{-1.5em}
\caption{\textbf{Teaser: Instruction-Conditioned Safety Alignment.} Same prompt with different alignment instructions. Only CITA produces valid, coherent responses in both modes.}
\label{fig:teaser_safety}
\vspace{-1.5em}
\end{figure*}

%Consider an enterprise agent platform in which the same backbone model is invoked under distinct workflow contracts:


%In such deployments, static alignment forces either a growing zoo of checkpoints (each with its own evaluation, monitoring, and legal sign-off) or a lowest-common-denominator policy that underserves most workflows. Instruction-based alignment offers a third path: \textbf{one checkpoint}, \textbf{multiple alignment postures}, activated by \textbf{explicit, auditable alignment instructions} inside the agent’s orchestration layer.

%\vspace{-1em}

\subsection{Introducing: ECLIPTICA \& CITA}
\label{sec:solution}

We present \textbf{ECLIPTICA}, a framework for \textbf{instruction-driven alignment} that elevates alignment to a \textbf{runtime-switchable interface} for agent platforms.
We instantiate ECLIPTICA with \textbf{CITA} (\textbf{C}ontrastive \textbf{I}nstruction-\textbf{T}uned \textbf{A}lignment), a training algorithm that teaches models to \textbf{bind} alignment instructions to preference behavior and to \textbf{remain stable under switching}.
CITA couples supervised fine-tuning with contrastive preference optimization and a \textbf{mandatory KL anchor} that prevents collapse and preserves a controllable policy geometry.
Figure~\ref{fig:teaser_safety} demonstrates the target behavior: for the \emph{same} user prompt, CITA supports \textbf{reliable switching} between strict refusal and permissive, harm-minimizing guidance, while standard trainers often exhibit repetition, leakage, incoherence, or unsafe drift.

\vspace{-0.35em}
\begin{enumerate}[leftmargin=*,itemsep=0em]
    %\item \textbf{ECLIPTICA Framework:} We formalize inference-time alignment control as a deployment primitive: \emph{alignment is an explicit instruction channel}, not a frozen property of weights.

     \item \textbf{ECLIPTICA benchmark:} We release a controlled diagnostic benchmark of \textbf{3,000} cases ($300$ prompts $\times$ $10$ instruction types) that isolates the causal effect of alignment instructions by holding the user prompt fixed.

    \item \textbf{CITA Algorithm:} We introduce an objective that combines SFT and contrastive preference learning with \textbf{mandatory KL regularization}.
    Unlike standard DPO where KL is often treated as optional, it is \textbf{structural} for \emph{switchability}.

    \item \textbf{Comprehensive Evaluation:} On \textbf{Llama-3.1-8B}~\cite{dubey2024llama}, across five deterministic benchmarks, \textbf{CITA achieves \colorbox{bg-green}{95.6\%} instruction-alignment efficiency}---outperforming DPO (70.5\%) by 25.1 pp, GRPO (59.3\%) by 36.3 pp, PPO (50.4\%) by 45.2 pp, and SFT (15.9\%) by 79.7 pp. On TruthfulQA~\cite{lin2022truthfulqa}, CITA shows \textbf{54$\times$ stronger adaptation} than DPO (+0.054 vs +0.001).
\end{enumerate}


\paragraph{ECLIPTICA vs.\ inference-time alignment methods.}
Recent work explores \emph{runtime controllability} via \textbf{attribute-conditioned steering} and \textbf{multi-objective decoding}
(e.g.,~\cite{dong-etal-2023-steerlm,wang-etal-2024-arithmetic,shi2024decodingtime}),
and via \textbf{inference-/decoding-time} value or reward guidance~\cite{liu-etal-2024-inference,chen2025pad}.
\textbf{ECLIPTICA} targets a distinct deployment primitive: \textbf{one checkpoint} that internalizes \emph{alignment instructions} as a \textbf{first-class control channel},
learning an \emph{instruction-conditioned preference policy} with \textbf{reliable, bidirectional switching}
(\emph{strict refusal} $\leftrightarrow$ \emph{permissive, harm-minimizing guidance}) \textbf{under a fixed user prompt}.
To our knowledge, these prior methods do not provide a \emph{preference-optimized}, instruction-conditioned objective \textbf{explicitly engineered} for this switching regime,
nor a controlled \emph{prompt-held-constant} benchmark (ISD) designed to \textbf{isolate the causal effect} of alignment instructions.
We provide a detailed comparison and design-differences table in Appendix~\ref{sec:appendix_related}.


\begin{comment}
\vspace{-2mm}
\begin{table}[H]
\centering
\footnotesize
\begin{tabularx}{\columnwidth}{@{}Xccccc@{}}
\toprule
\textbf{Property} & \textbf{SFT} & \textbf{PPO} & \textbf{GRPO} & \textbf{DPO} & \textbf{CITA} \\
\midrule
Reward Model Required & No & Yes & No & No & No \\
\hline
Preference Pairs & No & No & No & Yes & Yes \\
\hline
Instruction-Conditioned Switching & No & No & No & No & \textbf{Yes} \\
\hline
KL Regularization & --- & Yes & No & Implicit & \textbf{Yes} \\
\hline
Dynamic Policy & No & No & No & No & \textbf{Yes} \\
\hline
Agent-Compatible & No & No & No & No & \textbf{Yes} \\
\bottomrule
\end{tabularx}
\caption{Comparison of alignment methods. All trainers support instruction-augmented training (Instruct variants), but only CITA learns to dynamically switch behavior based on alignment instructions. PPO requires an external reward model; GRPO uses inline reward functions. DPO has implicit KL via $\beta$ temperature.}
\label{tab:method_comparison}
\vspace{-4mm}
\end{table}
\end{comment}

















